\SourceLeaf{432}%
\[
\begin{aligned}
Z&=\pm\frac{R}{\BareRoot{(PP+QQ+RR)}}\\
&=\pm\frac{1}{\BareRoot{(1+tt+uu)}},
\end{aligned}
\]
also
\[
AC-BB=\frac{TV-UU}{(1+tt+uu)^{2}}.
\]

\textls[100]{Jedem} Punkt der gekrümmten Fläche entspricht durch die dort gezogene Normale und den zu ihr parallelen Radius der Hilfskugel ein bestimmter Punkt $L$ auf der Oberfläche der Hilfskugel. Entsprechend bilden die Punkte auf der Hilfskugel, die allen Punkten einer \textls[100]{Linie} auf der gekrümmten Fläche entsprechen, eine Linie, die dieser Linie entspricht. Ebenso entspricht jeder begrenzten Figur auf der gekrümmten Fläche eine begrenzte Figur auf der Hilfskugel. Der Flächeninhalt der letzteren soll als Maß für die Amplitude der ersteren gelten. Wir werden diesen Flächeninhalt entweder als Zahl betrachten, wobei das Quadrat des Radius der Hilfskugel die Einheit ist, oder ihn in Grad usw. ausdrücken, indem wir den Flächeninhalt der Halbkugel $=360^{0}$ setzen.

Der Vergleich des Flächeninhalts auf der gekrümmten Fläche mit der entsprechenden Amplitude führt zu dem Begriff, den wir das Krümmungsmaß der Fläche nennen. Ist jener dieser proportional, so heißt die Krümmung gleichförmig. Der Quotient, den man durch Division der Amplitude durch den Flächeninhalt erhält, heißt dann Krümmungsmaß. Dies ist der Fall, wenn die gekrümmte Fläche eine Kugel ist. Das Krümmungsmaß ist dann ein Bruch mit dem Zähler $1$ und dem Quadrat des Radius als Nenner.

Wir betrachten das Krümmungsmaß als positiv, wenn die Grenzen die Figuren auf der gekrümmten Fläche und auf der Hilfskugel im gleichen Sinn umlaufen, und als negativ, wenn sie die Figuren im entgegengesetzten Sinn umlaufen. Besteht jene Proportionalität nicht, so ist die Fläche ungleichförmig gekrümmt. An jeder Stelle ergibt sich ein besonderes Krümmungsmaß aus dem Vergleich der einander entsprechenden unendlich kleinen Teile auf der gekrümmten Fläche und auf der Hilfskugel. Es sei $d\sigma$ ein Flächenelement auf jener, $d\Sigma$ das entsprechende auf der Hilfskugel. Dann ist
\[
\frac{d\Sigma}{d\sigma}
\]
das Krümmungsmaß an dieser Stelle.
